\appendix
\section{Storage layouts and constants}
\label{app:layout}

\subsection{Octree node record (host)}

\begin{center}
\begin{tabular}{@{}lll@{}}
\toprule
Field & Type & Bytes\\
\midrule
\code{sdf[8]} & float & 32\\
\code{brushIndex} & int & 4\\
\code{textureType} & int & 4\\
\code{data} & vec2 & 8\\
\code{hsv} & vec3 & 12\\
\code{blockId} & uint & 4\\
\code{bits} & uint8 & 1\\
\code{lod} & uint8 & 1\\
\code{chunkLod} & uint8 & 1\\
\code{version} & uint & 4\\
\bottomrule
\end{tabular}
\end{center}

\subsection{GPU node table (flattened)}

\begin{center}
\begin{tabular}{@{}lll@{}}
\toprule
Binding & Record & Bytes\\
\midrule
0 & \code{vec4 sdf0} & 16\\
1 & \code{vec4 sdf1} & 16\\
2 & \code{vec4 cubeMinSize} & 16\\
3 & \code{ivec4 meta} & 16\\
4 & \code{vec4 nodeHsv[2]} & 32\\
5 & \code{ivec4 child03} & 16\\
6 & \code{ivec4 child47} & 16\\
\midrule
total & & 128\\
\bottomrule
\end{tabular}
\end{center}

\subsection{Packed vertex record}

\begin{center}
\begin{tabular}{@{}lll@{}}
\toprule
Word & Field & Bytes\\
\midrule
0 & position.xyz + color.x & 16\\
1 & color.yz + texCoord.xy & 16\\
2 & normal.xyz + brush bits & 16\\
3 & textureType bits + hsv & 16\\
\midrule
total & & 64\\
\bottomrule
\end{tabular}
\end{center}

\subsection{Push constants}

\begin{center}
\begin{tabular}{@{}ll@{}}
\toprule
Field & Meaning\\
\midrule
\code{cellCount} & scan cells in this dispatch\\
\code{targetLod} & rung to scan\\
\code{fromId} & descent anchor\\
\code{fromBrush} & fallback brush\\
\code{rootMinSize} & world root $(x,y,z,L)$\\
\code{worldSize} & root edge for epsilon\\
\code{keyCapacity} & raw key bound\\
\code{pass} & 0 walk / 1 finalize\\
\code{vertexBase} & span-relative vertex base\\
\code{vertexCapacity} & vertex span bound\\
\code{indexBase} & span-relative index base\\
\code{indexCapacity} & index span bound\\
\code{indexCount} & raw vertices to remap\\
\code{nodeCount} & node-table entries to write\\
\code{diag} & diagnostic counters on/off\\
\bottomrule
\end{tabular}
\end{center}

\subsection{Constants}

\begin{center}
\begin{tabular}{@{}lll@{}}
\toprule
Name & Value & Purpose\\
\midrule
\code{kVerticesPerCellHint} & 25 & key sizing\\
\code{kKeyCapacityHeadroom} & 64 & key sizing\\
\code{kMaxKeysPerCall} & $2^{23}$ & hard key cap\\
\code{BREAK\_CAP} & 64 & final breaks per edge\\
\code{LOCAL\_BREAK\_CAP} & 32 & local breaks per interval\\
\code{TASK\_STACK} & 256 & interval stack\\
\code{POS\_MEMO} & 8 & QEF memo\\
\code{kInstanceStride} & 32 & vegetation instance\\
\code{kMaxVegChunks} & 4096 & vegetation chunk cap\\
\code{kPlacementScratchPerChunk} & 131072 & vegetation coarse cap\\
\bottomrule
\end{tabular}
\end{center}

\section{The complete scan-cell algorithm}
\label{app:algorithm}

\begin{algorithm}[H]
\caption{GPU scan cell (one invocation per cell)}
\begin{algorithmic}[1]
\State $cell \gets$ cell descriptor from table
\State $corners[8] \gets$ node SDF samples
\For{$e=0$ to $11$}
  \State $(i,j) \gets$ \code{EDGES}[e]
  \If{sign($corners[i]$) $=$ sign($corners[j]$)}
    \State \textbf{continue}
  \EndIf
  \State $span \gets$ \code{makeEdgeSpan}(e, cell)
  \If{$span.end-span.start \le 2\epsilon$}
    \State \textbf{continue}
  \EndIf
  \State $breaks \gets \{span.start, span.end\}$
  \State \code{collectBreaks}(span, span.start, span.end)
  \If{overflow} \State \textbf{fail chunk} \EndIf
  \State sort/unique($breaks$)
  \For{$k=1$ to $|breaks|-1$}
    \If{$breaks[k]-breaks[k-1] > 2\epsilon$}
      \State \code{emitSegment}(span, $breaks[k-1]$, $breaks[k]$,
        $sign_i$, $sign_j$)
    \EndIf
  \EndFor
\EndFor
\end{algorithmic}
\end{algorithm}

\section{Glossary}
\label{app:glossary}

\begin{description}[leftmargin=2.2em,style=nextline]
\item[Chunk] A cube of edge $\Delta$ whose subtree is flattened and meshed
  as one unit.  A chunk root is published as a draw entry.
\item[Rung] A level of the chunk ladder.  $r=0$ is the finest published
  level; larger $r$ are coarser ancestors.
\item[Scan cell] A cube at the target rung whose twelve edges are scanned
  for sign changes.
\item[Edge span] The interval of an active cube edge, with its fixed
  coordinates and epsilon.
\item[Break] A coordinate inside an edge span where a neighbouring cell
  boundary forces a split.
\item[Owner] The surface cell that determines the winding and leads the
  ring of a sub-segment.
\item[Raw key] The node index of one triangle corner, written in stage 0
  and remapped to a slot in stage 1.
\item[Slot] The dense vertex index assigned to a surface node by the host.
\item[Span] A reserved range in the packed vertex/index pools.
\item[Stamp] The per-entry LoD keep/reject decision written by the cull and
  consumed by the shadow/RT paths.
\end{description}
